% Einheit 3 von 4 – Klammern auflösen (bank/terme/e3.jsonl, zusammenbau v0.9)
\einheitenkopf[e3]{Einheit 3 von 4 · Klammern auflösen}
\zweigzeile{Hier lernst du, Klammern aufzulösen – Plusklammer, Minusklammer, Zahl mal Klammer (Distributivgesetz) · ab Kl. 7, je nach Buch bis Kl. 8 · keine P10-Aufgabe · baut auf: Addieren und Subtrahieren negativer Zahlen, Punkt vor Strich mit Zahlen}
\setcounter{aufgabe}{26}


% erkennung: terme-e3-k1-s0-v1, terme-e3-k1-s0-v2, terme-e3-k1-s0-v3, terme-e3-k1-s0-v4
\begin{kbaufgabe}{Ich erkenne, was vor einer Klammer steht.}
\begin{kbblock}
\kbanweisung{Was steht im Term … direkt vor der Klammer?}
\lbpaar{\kbteil{a)}{$7 + (a - 2)$ \quad \lbfeld}{}
}{\kbteil{b)}{$10 - (b + 4)$ \quad \lbfeld}{}}
\end{kbblock}
\begin{kbblock}
\lbpaar{\kbteil{c)}{$5 \cdot (x - 1)$ \quad \lbfeld}{}
}{\kbteil{d)}{$2y - 6 \cdot (y + 1)$ \quad \lbfeld}{}}
\end{kbblock}
\end{kbaufgabe}


% leiter: terme-e3-k2-s1-v1, terme-e3-k2-s1-v2, terme-e3-k2-s1-v3, terme-e3-k2-s2-v1, terme-e3-k2-s3-v1
\begin{kbaufgabe}{Ich kann Klammern auflösen.}
\begin{kbblock}
\kbanweisung{Löse die Klammer auf.}
\lbpaar{\kbteil{a)}{$5a + (2b - 1) =$ \lbfeld}{}
}{\kbteil{b)}{$5a + (2b - 3) =$ \lbfeld}{}}
\end{kbblock}
\begin{kbblock}
\lbpaar{\kbteil{c)}{$5a + (2b - 4) =$ \lbfeld}{}
}{\item[]}
\end{kbblock}
\begin{kbblock}
\kbanweisung{Löse die Klammer mit der Tabelle auf.}
\kbteil{d)}{$5 \cdot (x + 3)$}{}
\kbdarunter{\sachtabelle{ccc}{$\cdot$ & $x$ & $+3$}{$5$ & \leerzelle & \leerzelle}}
\kbantwort{= \kbfeld}
\end{kbblock}
\begin{kbblock}
\kbteil{e)}{Rechne $30 - (12 + 8)$ auf zwei Wegen: erst die Klammer ausrechnen und dann abziehen; dann Glied für Glied abziehen. Vergleiche beide Ergebnisse.}{}
\item[]\kbraum{2}
\end{kbblock}
\end{kbaufgabe}


% leiter: terme-e3-k2-s4-v1, terme-e3-k2-s5-v1, terme-e3-k2-s6-v1, terme-e3-k2-s7-v1
\begin{kbaufgabe}{Ich kann Klammern auflösen. – weiter}
\begin{kbblock}
\kbanweisung{Löse die Klammer auf.}
\lbpaar{\kbteil{a)}{$5a - (2b + 4) =$ \lbfeld}{}
}{\kbteil{b)}{$7a - (2b - 5 + c) =$ \lbfeld}{}}
\end{kbblock}
\begin{kbblock}
\lbpaar{\kbteil{c)}{$-2 \cdot (x + 5) =$ \lbfeld}{}
}{\item[]}
\end{kbblock}
\begin{kbblock}
\kbzweispaltig[0.46]{\kbspalte{\kbanweisung{Löse die Klammer auf und fasse zusammen.}\kbteil{d)}{$5x + 2 \cdot (x + 6)$}{}\kbantwort{= \kbfeld}}}{\kbraum{2}}
\end{kbblock}
\end{kbaufgabe}


% pruefung: terme-e3-k2-s8-v1
\begin{kbaufgabe}{Ich kann auch zwei Klammern in einem Term auflösen und zusammenfassen.}
\begin{kbblock}
\kbzweispaltig[0.46]{\kbspalte{\kbanweisung{Löse die Klammern auf und fasse zusammen.}\kbteil{}{$3 \cdot (2x + 5) - 2 \cdot (x - 4)$}{}\kbantwort{= \kbfeld}}}{\kbraum{2}}
\end{kbblock}
\end{kbaufgabe}


% pflicht: terme-e3-k3-s1-v1, terme-e3-k3-s1-v2, terme-e3-k3-s1-v3
\begin{kbaufgabe}{Ich finde den Fehler beim Auflösen von Klammern.}
\begin{kbblock}
\kbteil{a)}{Jan soll die Klammer in $14 - (a - 5)$ auflösen. Er rechnet so: \rechnung{14 - (a - 5) &= 14 - a - 5 &= 9 - a} Finde den Fehler und rechne richtig.}{}
\item[]\kbraum{2}
\end{kbblock}
\begin{kbblock}
\kbteil{b)}{Pia soll die Klammer in $9b - (4b - 2)$ auflösen. Sie rechnet so: \rechnung{9b - (4b - 2) &= 9b - 4b + 2 &= 5b + 2} Prüfe, ob Pia richtig gerechnet hat.}{}
\item[]\kbraum{2}
\end{kbblock}
\begin{kbblock}
\kbteil{c)}{Nils hat vier Klammern aufgelöst. Welche Ergebnisse können nicht stimmen? Begründe, ohne genau zu rechnen.}{}
\kbfrage{(1) $7 - (x + 2) = 5 - x$ \\ (2) $3 \cdot (y - 4) = 3y - 4$ \\ (3) $-(a - 6) = -a - 6$ \\ (4) $9 + (b - 1) = 8 + b$}
\item[]\kbraum{2}
\end{kbblock}
\end{kbaufgabe}


% pflicht: terme-e3-k3-s2-v1, terme-e3-k3-s2-v2, terme-e3-k3-s2-v3
\begin{kbaufgabe}{Ich kann begründen, wie man eine Klammer richtig auflöst.}
\begin{kbblock}
\kbteil{a)}{Begründe, ob die Terme $5 - (x - 2)$ und $7 - x$ gleichwertig sind.}{}
\item[]\kbraum{3}
\end{kbblock}
\begin{kbblock}
\kbteil{b)}{Kai sagt: „Bei $15 - (x - 3)$ muss ich nur das Vorzeichen von $x$ umdrehen. Das ergibt $15 - x - 3$.“ Begründe, ob Kai recht hat.}{}
\item[]\kbraum{3}
\end{kbblock}
\begin{kbblock}
\kbteil{c)}{Entscheide bei jeder Aussage, ob sie wahr oder falsch ist. Begründe, ohne genau zu rechnen.}{}
\kbfrage{(1) Steht ein Plus vor der Klammer, kann man die Klammer immer weglassen. \\ (2) Steht ein Minus vor der Klammer, bleibt das Vorzeichen des ersten Glieds in der Klammer immer gleich. \\ (3) Eine Zahl vor der Klammer wird mit jedem Glied in der Klammer malgenommen.}
\item[]\kbraum{3}
\end{kbblock}
\end{kbaufgabe}


% pflicht: terme-e3-k3-s3-v1, terme-e3-k3-s3-v2, terme-e3-k3-s3-v3
\begin{kbaufgabe}{Ich kann Klammern mit einer Tabelle oder einer Figur auflösen.}
\begin{kbblock}
\kbteil{a)}{Die Figur zeigt ein Rechteck. Die Maße sind in cm angegeben. Schreibe einen Term für den Flächeninhalt mit Klammer. Schreibe den Term dann ohne Klammer.}{}
\kbdarunter{\viereck[seiten={x+2,5,x+2,5}]{(0,0)}{(6,0)}{(6,2.5)}{(0,2.5)}}
\kbantwort{\kbfeld}
\end{kbblock}
\begin{kbblock}
\kbteil{b)}{Die Tabelle zeigt eine aufgelöste Klammer. Welcher Term mit Klammer gehört dazu?}{}
\kbdarunter{\sachtabelle{ccc}{$\cdot$ & $a$ & $+6$}{$2$ & $2a$ & $+12$}}
\kbantwort{\kbfeld}
\end{kbblock}
\begin{kbblock}
\kbteil{c)}{Ergänze die Tabelle. Schreibe dann den Term mit Klammer und den Term ohne Klammer auf.}{}
\kbdarunter{\sachtabelle{ccc}{$\cdot$ & $x$ & \leerzelle}{$4$ & \leerzelle & $+28$}}
\kbantwort{\kbfeld}
\end{kbblock}
\end{kbaufgabe}


% pflicht: terme-e3-k3-s4-v1, terme-e3-k3-s4-v2, terme-e3-k3-s4-v3
\begin{kbaufgabe}{Ich kann Terme mit Klammer im Alltag aufstellen und auflösen.}
\begin{kbblock}
\kbteil{a)}{Ein Kinoticket kostet 9 €. Eine Klasse mit 24 Schülern nimmt noch $x$ Begleiter mit. Schreibe die Gesamtkosten in Euro als Term mit Klammer. Löse die Klammer auf.}{}
\item[]\kbraum{2}
\end{kbblock}
\begin{kbblock}
\kbteil{b)}{Ein Rechteck ist $a$ cm lang und 12 cm breit. Schreibe den Umfang in cm als Term mit Klammer und löse die Klammer auf. Reicht ein Rahmenband von 60 cm, wenn $a = 20$ ist?}{}
\item[]\kbraum{2}
\end{kbblock}
\begin{kbblock}
\kbteil{c)}{Ein Muffin kostet 2,50 €. Für eine Feier werden 20 Muffins bestellt und noch $x$ Muffins dazu. Schreibe den Preis in Euro als Term mit Klammer. Löse die Klammer auf. Was kosten die Muffins bei $x = 4$?}{}
\item[]\kbraum{2}
\end{kbblock}
\end{kbaufgabe}
